% ARCHIVED v3.83: the corridor block at v3.82 (counts, S&C-only best-rule table, text-only frontier), replaced at the author's critique.
\subsection{UAV-corridor}\label{sec:res:uav:corridor}

\subsubsection{Generative Models}\label{sec:res:uav:corridor:raw}
The two constraints are read one at a time, as the avoiding geometries are: every table of this scene has a
\emph{tilt} block and a \emph{hump} block, and nothing is averaged across them.
Without projection every model reaches the end of the corridor, the scene's finish line
(\autoref{sec:setup:metrics:uav}), on every flight at every budget, and no flight is free of violations
under either constraint (\autoref{tab:uav-corridor-raw}): a flight violates the tilt on
$74.5$ to $99.4$ control steps on average and the hump on $28.6$ to $36.6$. The plans follow the level
lanes of the demonstrations, which were flown before either boundary existed
(\autoref{sec:setup:data:uav}). The flow-based models lie within one standard deviation of one another
under both constraints; the diffusion baseline has the most violating steps under the tilt and the fewest
over the hump. Because every configuration reaches the goal, no cost frontier is drawn before projection:
a frontier over points that share one outcome is a ranking by price (\autoref{sec:res:aligning:budget}),
and the price here is the budget's, from $9.1$ to $9.4$\,ms per action at one evaluation to $170.8$ to
$172.3$\,ms at twenty, and about $178$\,ms for the baseline.

\begin{table}[htpb]
  \caption[UAV-corridor: the generative models before projection]{UAV-corridor with the test-time
  constraints, \textbf{before projection}: the plan flown as generated, ten flights per row, spread over
  the three routes, one training seed, one block per constraint (\emph{tilt}, \emph{hump}). The grid is: the two average-velocity models at
  $\nfe\in\{1,2,3\}$, instantaneous-velocity matching at $\nfe\in\{1,2,3,5,20\}$, the diffusion baseline
  at its trained budget, printed last. \emph{Success}: crosses the finish line
  (\autoref{sec:setup:metrics:uav}). \emph{Violation-free}: no violating control step. \emph{S\&C}:
  success on a violation-free flight. \emph{Violating steps} and \emph{ms/step}: mean $\pm$ the sample
  standard deviation over the flights (\autoref{sec:setup:metrics:spread}). \baselinemark\ the
  diffusion model of \ac{DPCC}.}
  \label{tab:uav-corridor-raw}
  \centering
  \footnotesize
  \setlength{\tabcolsep}{3pt}
  \begin{tabularx}{\linewidth}{@{}L r rrr rr@{}}
    \toprule
    Model & $\nfe$ & Success & \shortstack{Violation-\\free} & S\&C & \shortstack{Violating\\steps} & ms/step \\
    \midrule
    \multicolumn{7}{@{}l}{\emph{tilt}}\\
    MeanFM & 1 & 10/10 & 0/10 & 0/10 & $91.3 \pm 34.8$ & $9.4 \pm 0.4$ \\
    MeanFM & 2 & 10/10 & 0/10 & 0/10 & $92.9 \pm 30.2$ & $18.1 \pm 0.4$ \\
    MeanFM & 3 & 10/10 & 0/10 & 0/10 & $92.0 \pm 30.4$ & $27.0 \pm 0.4$ \\
    CI-MeanFM, $\alpha_{\mathrm{end}}=0.2$ & 1 & 10/10 & 0/10 & 0/10 & $90.0 \pm 32.6$ & $9.3 \pm 0.4$ \\
    CI-MeanFM, $\alpha_{\mathrm{end}}=0.2$ & 2 & 10/10 & 0/10 & 0/10 & $93.8 \pm 30.0$ & $18.4 \pm 0.4$ \\
    CI-MeanFM, $\alpha_{\mathrm{end}}=0.2$ & 3 & 10/10 & 0/10 & 0/10 & $93.1 \pm 30.1$ & $27.3 \pm 0.4$ \\
    FM & 1 & 10/10 & 0/10 & 0/10 & $74.5 \pm 25.6$ & $9.1 \pm 0.4$ \\
    FM & 2 & 10/10 & 0/10 & 0/10 & $83.9 \pm 28.4$ & $17.7 \pm 0.4$ \\
    FM & 3 & 10/10 & 0/10 & 0/10 & $87.0 \pm 28.3$ & $25.9 \pm 0.5$ \\
    FM & 5 & 10/10 & 0/10 & 0/10 & $89.9 \pm 28.4$ & $43.4 \pm 0.4$ \\
    FM & 20 & 10/10 & 0/10 & 0/10 & $94.7 \pm 27.0$ & $170.8 \pm 0.4$ \\
    \mainconfig{\baselinemark\ Diffusion} & \mainconfig{20} & \mainconfig{10/10} & \mainconfig{0/10} & \mainconfig{0/10} & \mainconfig{$99.4 \pm 27.0$} & \mainconfig{$178.1 \pm 0.8$} \\
    \midrule
    \multicolumn{7}{@{}l}{\emph{hump}}\\
    MeanFM & 1 & 10/10 & 0/10 & 0/10 & $36.0 \pm 8.9$ & $9.4 \pm 0.4$ \\
    MeanFM & 2 & 10/10 & 0/10 & 0/10 & $35.6 \pm 8.1$ & $18.2 \pm 0.4$ \\
    MeanFM & 3 & 10/10 & 0/10 & 0/10 & $35.7 \pm 7.8$ & $27.2 \pm 0.4$ \\
    CI-MeanFM, $\alpha_{\mathrm{end}}=0.2$ & 1 & 10/10 & 0/10 & 0/10 & $36.6 \pm 8.4$ & $9.3 \pm 0.4$ \\
    CI-MeanFM, $\alpha_{\mathrm{end}}=0.2$ & 2 & 10/10 & 0/10 & 0/10 & $35.6 \pm 7.9$ & $18.0 \pm 0.4$ \\
    CI-MeanFM, $\alpha_{\mathrm{end}}=0.2$ & 3 & 10/10 & 0/10 & 0/10 & $35.6 \pm 7.6$ & $27.1 \pm 0.4$ \\
    FM & 1 & 10/10 & 0/10 & 0/10 & $34.7 \pm 6.8$ & $9.1 \pm 0.4$ \\
    FM & 2 & 10/10 & 0/10 & 0/10 & $32.6 \pm 7.8$ & $17.5 \pm 0.3$ \\
    FM & 3 & 10/10 & 0/10 & 0/10 & $32.6 \pm 8.0$ & $26.3 \pm 0.4$ \\
    FM & 5 & 10/10 & 0/10 & 0/10 & $32.0 \pm 7.9$ & $43.6 \pm 0.4$ \\
    FM & 20 & 10/10 & 0/10 & 0/10 & $31.3 \pm 7.6$ & $172.3 \pm 2.6$ \\
    \mainconfig{\baselinemark\ Diffusion} & \mainconfig{20} & \mainconfig{10/10} & \mainconfig{0/10} & \mainconfig{0/10} & \mainconfig{$28.6 \pm 7.7$} & \mainconfig{$178.2 \pm 2.7$} \\
    \bottomrule
  \end{tabularx}
\end{table}
\dataref{R33 (\texttt{PENDING\_20260923\_uav\_corridor\_v3\_pillars\_v2\_paper\_runs.md}), landed 24-09:
\texttt{temp/23-09-FULL/24-09-1000/batch\_uav\_20260924\_081422/per\_rollout\_detail.csv}, tags
\texttt{p23cv3t} (tilt) and \texttt{p23cv3ah} (hump), CI-MeanFM \texttt{EPlatest\_p23cv3\{t,ah\}};
geometries \texttt{corridor\_v3\_tilt}, \texttt{corridor\_v3\_ablation\_hump}; variant \texttt{diffuser};
the first ten flights of each cell (routes L, C, R cycling, 4/3/3); columns \texttt{success\_relaxed},
\texttt{collision\_free\_completed}, \texttt{n\_success\_relaxed\_and\_constraints}, \texttt{n\_violations},
\texttt{avg\_time\_ms}. Analysis of record \texttt{DA\_in\_Paper/analysis/DA\_20260924\_corridor\_v3.md}
\S3.1 (\texttt{corridor\_v3\_grid.py}); every cell recomputed from the CSV at v3.80 and equal. The v2
corpus this table was computed from until v3.65 is archived.}

\FloatBarrier
\subsubsection{Projection}\label{sec:res:uav:projected}
Under the tilt only the diffusion baseline succeeds violation-free (\autoref{tab:uav-corridor}). Under
per-step projection every one of its flights reaches the end of the corridor, and eight of ten do so
without a violation under cumulative-cost selection, three under each of the other two rules, at $651$ to
$680$\,ms per action. No flow-based configuration does, at any budget and under either projection method:
every projected flow-based flight reaches the end of the corridor, and every one violates the leaned plane
on a few control steps just before it. Over the hump violation-free successes come from two
configurations: the baseline under per-step projection, on ten flights of ten under random selection at
$654$\,ms per action and nine under the other two rules, and instantaneous-velocity matching at twenty
evaluations under endpoint projection, best with a single candidate, on eight at $291$\,ms. No flow-based
model succeeds over the hump under per-step projection: at one and two evaluations the flights stop short
of the apex, and from three evaluations on every flight crosses it with a few violating steps at the apex.
At the route's goal point, the line these flights were first scored at, the baseline's counts are $0/10$
under the tilt and $3$ to $6$ of ten over the hump, and every flow-based count is the same
(\autoref{sec:setup:metrics:uav}): $41$ of the baseline's $60$ projected flights had passed the end of the
corridor and were still flying towards the goal when the episode limit ended them. The budgets are read
as one ladder across the scene: the per-step projector is active at every budget, endpoint projection at
$\nfe\ge 3$, where it has a guiding step (\autoref{sec:res:constraints:degenerate}), so at $\nfe=1$ and
$\nfe=2$ the table carries per-step rows only.

\begin{table}[htpb]
  \caption[UAV-corridor: the generative models after projection]{UAV-corridor, one block per test-time
  constraint, \textbf{after projection}: each model at each budget under the per-step projector of \ac{DPCC}
  (activation threshold $0.5$, tightened constraints) and, from $\nfe=3$, under endpoint projection, each at
  its best selection rule (most successes with constraint satisfaction, then fewest violating steps, then
  least time; rules as in \autoref{tab:uav-corridor-projection}). Ten flights per cell, one training seed.
  \emph{S\&C}: success on a violation-free flight; \emph{ms/step}: time to compute one action; \emph{---}:
  no guiding step. At the route's goal point the baseline reads $0/10$ (tilt) and $6/10$ under $t$ (hump);
  every other row is unchanged. \baselinemark\ the diffusion model of \ac{DPCC}.}
  \label{tab:uav-corridor}
  \centering
  \footnotesize
  \setlength{\tabcolsep}{4pt}
  \begin{tabularx}{\linewidth}{@{}L r rrr rrr@{}}
    \toprule
    & & \multicolumn{3}{c}{per-step projection} & \multicolumn{3}{c}{endpoint projection} \\
    \cmidrule(lr){3-5}\cmidrule(l){6-8}
    Model & $\nfe$ & S\&C & rule & ms/step & S\&C & rule & ms/step \\
    \midrule
    \multicolumn{8}{@{}l}{\emph{tilt}}\\
    MeanFM & 1 & 0/10 & $c$ & 27.4 & --- & --- & --- \\
    MeanFM & 2 & 0/10 & $c$ & 36.4 & --- & --- & --- \\
    MeanFM & 3 & 0/10 & $c$ & 93.6 & 0/10 & $c$ & 73.5 \\
    CI-MeanFM, $\alpha_{\mathrm{end}}=0.2$ & 1 & 0/10 & $c$ & 27.9 & --- & --- & --- \\
    CI-MeanFM, $\alpha_{\mathrm{end}}=0.2$ & 2 & 0/10 & $c$ & 36.1 & --- & --- & --- \\
    CI-MeanFM, $\alpha_{\mathrm{end}}=0.2$ & 3 & 0/10 & $c$ & 93.4 & 0/10 & $c$ & 72.7 \\
    FM & 1 & 0/10 & $c$ & 27.6 & --- & --- & --- \\
    FM & 2 & 0/10 & $c$ & 34.7 & --- & --- & --- \\
    FM & 3 & 0/10 & $c$ & 79.2 & 0/10 & $c$ & 71.6 \\
    FM & 5 & 0/10 & $c$ & 116.9 & 0/10 & $c$ & 115.2 \\
    FM & 20 & 0/10 & $c$ & 360.4 & 0/10 & $c$ & 416.9 \\
    \mainconfig{\baselinemark\ Diffusion} & \mainconfig{20} & \mainconfig{8/10} & \mainconfig{$c$} & \mainconfig{666.3} & --- & --- & --- \\
    \midrule
    \multicolumn{8}{@{}l}{\emph{hump}}\\
    MeanFM & 1 & 0/10 & $c$ & 26.4 & --- & --- & --- \\
    MeanFM & 2 & 0/10 & $r$ & 35.0 & --- & --- & --- \\
    MeanFM & 3 & 0/10 & $c$ & 91.1 & 0/10 & $c$ & 72.1 \\
    CI-MeanFM, $\alpha_{\mathrm{end}}=0.2$ & 1 & 0/10 & $t$ & 26.6 & --- & --- & --- \\
    CI-MeanFM, $\alpha_{\mathrm{end}}=0.2$ & 2 & 0/10 & $r$ & 35.4 & --- & --- & --- \\
    CI-MeanFM, $\alpha_{\mathrm{end}}=0.2$ & 3 & 0/10 & $c$ & 89.2 & 0/10 & $c$ & 71.9 \\
    FM & 1 & 0/10 & $c$ & 26.2 & --- & --- & --- \\
    FM & 2 & 0/10 & $r$ & 34.5 & --- & --- & --- \\
    FM & 3 & 0/10 & $c$ & 77.1 & 0/10 & $c$ & 71.9 \\
    FM & 5 & 0/10 & $c$ & 114.1 & 0/10 & $c$ & 112.4 \\
    FM & 20 & 0/10 & $c$ & 355.2 & 8/10 & single & 291.0 \\
    \mainconfig{\baselinemark\ Diffusion} & \mainconfig{20} & \mainconfig{10/10} & \mainconfig{$r$} & \mainconfig{654.2} & --- & --- & --- \\
    \bottomrule
  \end{tabularx}
\end{table}
\dataref{R33, the batch and tags of \autoref{tab:uav-corridor-raw}; per-step = \texttt{dpcc-\{r,c,t\}}
at $\eta=0.5$, endpoint = \texttt{hardflow\_sls\{,-r,-c,-t\}}, both with the corridor's stack
\texttt{-bounds\_free-pdes-tightened} (\autoref{sec:setup:protocol:uav}). Best rule = most
\texttt{n\_success\_relaxed\_and\_constraints}, then the lowest mean \texttt{n\_violations}, then the lowest
\texttt{avg\_time\_ms} (DA \S3.2); ties on S\&C are almost everywhere, since S\&C is 0. Every cell
recomputed from the CSV at v3.80 and equal. v3.82: success at the end of the corridor, $x=2.0$\,m (R45a,
\texttt{DA\_20260924\_corridor\_v3\_R45a\_clear\_line.md} \S2.1, read after each step as the evaluation
latches); only the baseline row changes (tilt $0\to8$, $c$; hump $6$ $t\to10$ $r$). On recorded start
positions only, hump $r$ reads $9/10$ and the best rule would be $c$ ($9/10$, $623.1$\,ms): one flight is
$1.2$\,mm short before its last step and $16.7$\,mm past after it (DA \S3). The unprojected twins are
\autoref{tab:uav-corridor-raw}.}

The budget makes the violations smaller and does not remove them. Wherever the flights cross the line,
the violating steps per flight fall with the budget, under both constraints and both projection methods.
Under the tilt they fall from
$9.8$ to $11.1$ at one and two evaluations to $2.1$ to $3.2$ at twenty under per-step projection and to
$1.4$ to $2.5$ under endpoint projection, and the deepest of them from $8.8$ to $11.9$\,cm to $1.3$ to
$2.2$ and $1.0$ to $1.3$\,cm. Over the hump they fall from $4.6$ to $6.5$ under per-step and $8.0$ to
$10.4$ under endpoint projection at three evaluations to $1.8$ to $2.2$ and $0.2$ to $0.7$ at twenty.
Under the tilt no budget reaches zero; over the hump only endpoint projection at twenty evaluations does,
on four to eight flights of ten, and it is the one flow-based configuration that succeeds violation-free.
The budget at
which a flight stops reaching the finish line lies at the other end of the ladder: under per-step
projection at one and two evaluations $177$ of the $180$ flow-based flights over the hump stop short of
the apex, and the other three reach the end of the corridor only after drifting across the boundary of its
wall. The baseline, at its one budget, violates the hump on $0.0$ to $0.4$ control steps per flight under
per-step projection and is the slowest planner on the scene: its projected flights leave the corridor a
median of $29$ and $39$ control steps after the flow-based ones (tilt, hump) and then approach the goal
point slowly and from the side, at $623$ to $680$\,ms per action, the highest price of the scene.

The frontier of this arm is short enough to state. Under the tilt it is one configuration, the baseline
under per-step projection with cumulative-cost selection: eight violation-free successes at $666$\,ms per
action. Over the hump it has two points and neither dominates the other: instantaneous-velocity matching at
twenty evaluations under endpoint projection with a single candidate, eight at $291$\,ms, and the baseline
under per-step projection with random selection, ten at $654$\,ms --- two more violation-free successes
for $2.2$ times the time per action.

The strict count hides what separates the flow-based planners from the baseline on this scene. A projected
flow-based flight that violates the leaned plane does so on a few control steps at its end --- about ten
at one and two evaluations, two or three at twenty --- and is violation-free on the rest. Read with the
near-satisfaction bar of \autoref{sec:setup:metrics:uav}, which counts a flight that succeeds with at least
$95\,\%$ of its control steps free of violations, every projected configuration reaches ten flights of
ten, except per-step projection at one and two evaluations over the hump, whose flights stop short of the
apex. The frontier is then the cheapest configuration: under the tilt one evaluation with per-step
projection, $27.4$ to $27.9$\,ms per action for the three flow-based models, and over the hump three
evaluations with endpoint projection, $71.9$ to $72.1$\,ms --- a twenty-fourth and a ninth of the
baseline's time per action at the same count. The three flow-based models tie at both points, within two
per cent of one another. The bar decides this reading: at $97\,\%$ the one- and two-evaluation cells under
the tilt fall to at most four flights of ten and the tilt's frontier becomes a staircase that ends at
instantaneous-velocity matching with three evaluations under per-step projection, ten flights at
$79$\,ms, while over the hump it stays at three evaluations with endpoint projection; below about
$87\,\%$ the bar would count unprojected flights over the hump as satisfying.
\dataref{Strict frontier at the end of the corridor: \texttt{DA\_20260924\_corridor\_v3\_R45a\_clear\_line.md}
\S2.2 (tilt: the baseline $c$ is the only eligible point; hump: two non-dominated points). Near-satisfaction
frontier: the same DA, Appendix C (every configuration strict / $\ge0.95$ / $\ge0.97$; $0.90$ gives the
same frontier; the floor, unprojected flights $\approx66\,\%$ tilt and $87\,\%$ hump violation-free; the
bar forgives up to $\approx14$ steps, $\approx0.4$\,s, up to $10$\,cm deep at $\nfe\le2$). Under the old
goal-plane line (DA of record \S3.4) the tilt had no eligible point and the hump one, FM $\nfe=20$
endpoint single. Frontier rule of
\texttt{corridor\_v3\_grid.frontier} (costs within one per cent are one cost). The builder
\texttt{frontier.py::fig\_uav\_corridor\_tradeoff} still draws the archived v2 corpus and is not used
(v3.80). Budget facts: DA \S3.2, recomputed from the CSV; depths from the flown paths, DA Appendix A.}

Where the violating steps lie is seen in the flights themselves (\autoref{fig:uav-corridor-side}), drawn
for instantaneous-velocity matching, the one flow-based model flown at twenty evaluations, without
projection and under per-step projection at one and at twenty evaluations. Under the tilt the unprojected
flights stay level, so their margin to the leaned plane shrinks steadily along the corridor, to as much
as $0.37$\,m inside the plane at its end. The projected flights move sideways and down until they hold a few
centimetres clear of the plane, and every one of them cuts into it in the last $0.3$\,m of the plane's
span: every violating step of every projected flight, on every model, budget and projection method, lies
between $x=1.71$ and $2.00$\,m. At the first violating step of each projected flow-based flight the
commanded position is already past the end of the plane, at $x=2.27$ to $2.52$\,m and $0.49$ to
$0.61$\,m ahead of the vehicle, and at every violating step the commanded position itself is clear of the
constraint. The plane is enforced over its span of $x$, and on this scene the span is checked at the
commanded position: once that position has left the span, the plan is released from the plane and
returns toward its route, and the vehicle, half a metre behind, follows it into the plane. The violating
steps under the tilt are thus the vehicle catching up with a plan that has already left the constraint,
which is why the budget makes them shallower without removing them. The baseline's projected flights leave
the corridor further to the side of their routes --- a median of $0.37$\,m, against $0.19$\,m for the
flow-based flights --- and stay there, and three to eight of ten of them fly the scene without a
violation.

Over the hump the unprojected flights stay level and pass through the roof. Under per-step projection at one
evaluation every flight climbs along the boundary and stops $0.12$ to $0.13$\,m short of the apex, at
$1.42$ to $1.43$\,m, where it stays until its $396$ control steps run out. At one and two evaluations
all three flow-based models stop this way, on $174$ of their $180$ flights, $0.01$ to $0.17$\,m short of
the apex; the other six clear the apex and then drift sideways across the boundary of the corridor wall.
Two properties of the setup meet at the apex. The plans never rise above $1.30$\,m, the top of the
altitude range of the demonstrations (dotted in \autoref{fig:uav-corridor-side}), while the apex asks for
$1.48$\,m. And each face of the roof, like the tilt, is enforced over its span of $x$ checked at the
commanded position, so while that position is short of the apex the rising face binds the whole plan,
including its part beyond the apex, where the face keeps rising. Both hold at every budget. At one and
two evaluations the corrected plan settles short of the apex; from three on it does not, and at twenty
the flights cross the apex grazing the boundary by about a centimetre and descend toward their route
altitude before the line.

\begin{figure}[htpb]
  \centering
  \includegraphics[width=\linewidth]{fig_uav_corridor_side}
  \caption[UAV-corridor: the flights along the corridor]{UAV-corridor along the corridor,
  instantaneous-velocity matching, ten flights per series, one training seed: without projection at
  $\nfe=1$ (grey) and under per-step projection with cumulative-cost selection at $\nfe=1$ and $\nfe=20$
  (blue). (a, b) Tilt: the margin of the flown position to the leaned plane, less the rotor reach, over
  the plane's span of $x$; a control step below the dashed line violates. (c, d) Hump: altitude over $x$,
  the roof shaded, the violation boundary dashed, the highest planned altitude dotted; a circle ends a
  flight short of the finish line.}
  \label{fig:uav-corridor-side}
\end{figure}
\dataref{Store figure \texttt{DA\_in\_Paper/figures/da/fig\_uav\_corridor\_side} (builder
\texttt{builders/altitude.py::fig\_uav\_corridor\_side}; extract \texttt{extract/corridor\_v3\_side.py}
$\to$ \texttt{data/corridor\_v3\_side.json}), cells read through \texttt{analysis/corridor\_v3\_grid.py}:
FM, \texttt{diffuser} at $\nfe=1$ and \texttt{dpcc-c-bounds\_free-pdes-tightened} at 1 and 20, the first
ten flights; $c$ is FM's best per-step rule at both budgets and both constraints (\autoref{tab:uav-corridor}).
Margin $=\mathbf{n}\cdot(\mathbf{p}-\mathbf{p}_0)-0.31$ with the evaluation's own \texttt{\_hs\_normal3};
hump boundary $=$ roof $+\,0.31\sqrt{1+s^2}$, $s=1.10/1.5$; highest planned altitude $1.2992$\,m, the
largest planned $z$ of the drawn cells. Violating-step positions of all $560$ projected tilt flights,
$4115$ steps: $1.709\le x\le 2.000$, $99\,\%$ at $x\ge 1.75$ (the DA's $[1.75,2.00]$ is its 5--95\,\%
range). Commanded position at the first violating step, $530$ projected flow-based flights: $x=2.27$--$2.52$
(median $2.38$), $0.49$--$0.61$\,m ahead (median $0.55$); clear of the constraint at $4014$ of $4014$
violating steps. Span check at the commanded position under \texttt{-pdes}:
\texttt{eval\_mix\_uav.py:1830--1833} (\texttt{\_gate\_x = p\_des[0]}) and \texttt{:1483--1484}. Hump
stops: per flight $x=-0.174$ to $-0.006$, $z=1.38$ to $1.51$\,m, every flight on the $396$-step cap; the six
others are MeanFM and CI-MeanFM at $\nfe=2$ under $t$, three each, wall violations with no collision
(\texttt{phys\_contact\_frac} $0$), ending $0.95$--$1.79$\,m from the goal. Plan bound: every stored plan
$z\in[0.9008,1.2992]$\,m (DA \S4.2, \texttt{LimitsNormalizer.unnormalize}); a third cause the DA names ---
the correction is the smallest in normalised coordinates --- is not stated in the text. Replaces the
three placeholders \texttt{fig:uav-corridor-tradeoff}, \texttt{-paths}, \texttt{-altitude} (v3.80).}

\paragraph{Altitude.} Both constraints of this scene act on altitude, and that is what they are for: under
the tilt the room is below the launch altitude, so a projected flight has to descend as it moves sideways;
under the hump it has to climb. Unprojected, every model flies level, between $0.83$ and $1.27$\,m where
the constraints act. Projected, the lowest altitude under the tilt drops on $540$ of the $560$ projected
flights, by a median of $0.12$ to $0.37$\,m per configuration, and the highest altitude over the hump
rises on all $560$, by a median of $0.27$ to $0.39$\,m and on single flights by up to $0.60$\,m
(\autoref{tab:uav-corridor-altitude}). Under the tilt the descent deepens with the budget, from $0.12$ to
$0.19$\,m for the two average-velocity models at one and two evaluations to $0.21$ to $0.26$\,m at three,
and from $0.19$ to $0.25$\,m for instantaneous-velocity matching at one and two to $0.35$ to $0.37$\,m at
twenty; the climb over the hump does not.

\begin{table}[htpb]
  \caption[UAV-corridor: the altitude change under projection]{UAV-corridor, the altitude change that
  projection makes: under the tilt the lowest altitude over $0.5\le x\le 2.0$\,m, over the hump the
  highest over $-0.5\le x\le 0.5$\,m, each the projected flight minus the unprojected flight of the same
  trial, as the median over a cell's ten flights. The range runs over the model's projected cells: every
  budget, projection method and selection rule. One training seed. \baselinemark\ the diffusion model of
  \ac{DPCC}.}
  \label{tab:uav-corridor-altitude}
  \centering
  \footnotesize
  \setlength{\tabcolsep}{5pt}
  \begin{tabularx}{\linewidth}{@{}L l rr@{}}
    \toprule
    Model & $\nfe$ & \shortstack{tilt: lowest altitude,\\median change (m)} & \shortstack{hump: highest altitude,\\median change (m)} \\
    \midrule
    MeanFM & 1, 2, 3 & $-0.26$ to $-0.13$ & $+0.33$ to $+0.39$ \\
    CI-MeanFM, $\alpha_{\mathrm{end}}=0.2$ & 1, 2, 3 & $-0.26$ to $-0.12$ & $+0.33$ to $+0.39$ \\
    FM & 1, 2, 3, 5, 20 & $-0.37$ to $-0.19$ & $+0.27$ to $+0.36$ \\
    \mainconfig{\baselinemark\ Diffusion} & \mainconfig{20} & \mainconfig{$-0.30$ to $-0.22$} & \mainconfig{$+0.32$ to $+0.34$} \\
    \bottomrule
  \end{tabularx}
\end{table}
\dataref{DA \S3.6 and Appendix B (median per cell and rule), recomputed from the npz at v3.80 with
\texttt{corridor\_v3\_grid.altitude}; flights paired by trial index. $540$ of $560$ tilt flights descend
(the $20$ that do not: MeanFM $8$, CI-MeanFM $10$, FM $2$, each within $+0.11$\,m); $560$ of $560$ hump
flights climb. The unprojected band is the range of the same window extremes. The v3.65 pilots
($-0.15$ to $-0.29$, $+0.30$ to $+0.40$) lie inside these ranges.}

None of the three quadrotor scenes is drawn as a cost frontier. On UAV-corridor it has at most two
points under either reading and is stated above.
UAV-pillars' result is the difference between a plan's outcome on D3IL-avoiding and the outcome of the
same plan flown, which is a comparison and not a trade-off (\autoref{sec:res:uav:models}). UAV-s-curve
has data, but no configuration flown on it succeeds with constraint satisfaction
(\autoref{tab:uav-scurve}); a cost frontier over points that all share one outcome is a ranking of
failures by price, so that scene gets the abort count printed beside each of its rows instead.

\FloatBarrier
\subsubsection{Projection Methods}\label{sec:res:uav:projection}
Per-step and endpoint projection separate on this scene only at twenty evaluations
(\autoref{tab:uav-corridor-projection}). At three and five evaluations the two are level, at no
violation-free success under either constraint; endpoint projection is as cheap or cheaper, $44$ to
$74$\,ms per action against $77$ to $94$ at three evaluations and $73$ to $115$ against $112$ to $118$ at
five, and it leaves more violating steps: at three evaluations about two more per flight under the tilt
($6.9$ to $9.2$ against $5.0$ to $7.5$) and about four more over the hump ($8.0$ to $10.4$ against $4.6$
to $6.5$), at five about one more. At twenty evaluations endpoint projection leaves fewer and shallower
violating steps than per-step projection under the tilt, $1.4$ to $2.5$ per flight and at most
$1.3$\,cm deep against $2.1$ to $3.2$ and $2.2$\,cm, and
over the hump it is the only flow-based configuration that succeeds violation-free, on four to eight
flights of ten depending on the rule, where per-step projection succeeds on none. The bar this thesis set
for endpoint projection, to overtake the per-step projector of \ac{DPCC} at a lower cost, is met on one
constraint at one budget: over the hump at twenty evaluations with a single candidate, eight
violation-free successes against none, at $291$ against $355$\,ms per action. With four candidates it
still overtakes the per-step projector there, at a higher cost ($422$ to $425$\,ms).

\begin{table}[htpb]
  \caption[UAV-corridor: the two projection methods under every selection rule]{UAV-corridor, success with
  constraint satisfaction --- success on a violation-free flight --- for both projection methods under
  every selection rule, one block per constraint (\emph{tilt}, \emph{hump}), at the budgets at which
  endpoint projection has a guiding step; ten flights
  per cell, tightened constraints, activation threshold $0.5$, one training seed. Per-step: $r$ random,
  $c$ cumulative projection cost, $t$ temporal consistency, four candidate plans. Endpoint:
  \emph{single} one candidate, then the three rules over four. The diffusion baseline has per-step rows
  only; at the route's goal point it reads $0/10$ under every rule (tilt) and $3$, $3$, $6$ (hump).}
  \label{tab:uav-corridor-projection}
  \centering
  \footnotesize
  \setlength{\tabcolsep}{5pt}
  \begin{tabularx}{\linewidth}{@{}L r rrr rrrr@{}}
    \toprule
    & & \multicolumn{3}{c}{Per-step} & \multicolumn{4}{c}{Endpoint} \\
    \cmidrule(lr){3-5}\cmidrule(l){6-9}
    Model & $\nfe$ & $r$ & $c$ & $t$ & single & $r$ & $c$ & $t$ \\
    \midrule
    \multicolumn{9}{@{}l}{\emph{tilt}}\\
    MeanFM & 3 & 0/10 & 0/10 & 0/10 & 0/10 & 0/10 & 0/10 & 0/10 \\
    CI-MeanFM, $\alpha_{\mathrm{end}}=0.2$ & 3 & 0/10 & 0/10 & 0/10 & 0/10 & 0/10 & 0/10 & 0/10 \\
    FM & 3 & 0/10 & 0/10 & 0/10 & 0/10 & 0/10 & 0/10 & 0/10 \\
       & 5 & 0/10 & 0/10 & 0/10 & 0/10 & 0/10 & 0/10 & 0/10 \\
       & 20 & 0/10 & 0/10 & 0/10 & 0/10 & 0/10 & 0/10 & 0/10 \\
    \mainconfig{\baselinemark\ Diffusion} & \mainconfig{20} & \mainconfig{3/10} & \mainconfig{8/10} & \mainconfig{3/10} & --- & --- & --- & --- \\
    \midrule
    \multicolumn{9}{@{}l}{\emph{hump}}\\
    MeanFM & 3 & 0/10 & 0/10 & 0/10 & 0/10 & 0/10 & 0/10 & 0/10 \\
    CI-MeanFM, $\alpha_{\mathrm{end}}=0.2$ & 3 & 0/10 & 0/10 & 0/10 & 0/10 & 0/10 & 0/10 & 0/10 \\
    FM & 3 & 0/10 & 0/10 & 0/10 & 0/10 & 0/10 & 0/10 & 0/10 \\
       & 5 & 0/10 & 0/10 & 0/10 & 0/10 & 0/10 & 0/10 & 0/10 \\
       & 20 & 0/10 & 0/10 & 0/10 & 8/10 & 6/10 & 6/10 & 4/10 \\
    \mainconfig{\baselinemark\ Diffusion} & \mainconfig{20} & \mainconfig{10/10} & \mainconfig{9/10} & \mainconfig{9/10} & --- & --- & --- & --- \\
    \bottomrule
  \end{tabularx}
\end{table}
\dataref{R33, the batch and tags of \autoref{tab:uav-corridor-raw}: variants \texttt{dpcc-\{r,c,t\}} and
\texttt{hardflow\_sls\{,-r,-c,-t\}} on the corridor projection stack;
\texttt{n\_success\_relaxed\_and\_constraints}; DA \S3.3, recomputed from the CSV at v3.80 and equal; v3.82: the
baseline row at the end of the corridor (R45a DA \S2, tilt $0/0/0\to3/8/3$, hump $3/3/6\to10/9/9$), every
flow-based row unchanged. The
DA also has the $\nfe=1,2$ per-step rows: all $0/10$. Mean violating steps and ms per rule (DA \S3.2, the
CSV): at $\nfe=3$ per-step $5.0$--$7.5$ (tilt) and $4.6$--$6.5$ (hump) at $77$--$94$\,ms, endpoint
$6.9$--$9.2$ and $8.0$--$10.4$ at $44$--$74$\,ms; at $\nfe=20$ FM per-step $354$--$360$\,ms, endpoint
single $289$--$291$, four candidates $417$--$425$. The DA's \enquote{within a projector the rules differ by at
most two violating steps} does not hold everywhere (endpoint, hump, CI-MeanFM $\nfe=3$: $8.0$ to $10.4$;
the baseline under the tilt: $0.8$ to $6.1$) and is not used.}
\guard{Endpoint projection has no row at $\nfe=1$ or $\nfe=2$ on this scene and will not get one: at the
activation threshold of $0.5$ the only active sampling step there is the terminal one, so there is no
guiding step for the method to act on (\autoref{sec:res:constraints:degenerate}). That is a property of
the method at those budgets, not a pending evaluation. The baseline's dashes are the opposite case:
endpoint projection has never been run on the diffusion model on any scene.}

The projection configuration of this scene differs from that of the other quadrotor scene
(\autoref{sec:setup:protocol:uav}), so the scenes are not pooled.

\FloatBarrier
